\documentclass[10pt,a4paper]{amsart}
\usepackage[margin=19mm]{geometry}
\usepackage{amsmath,amssymb,mathtools}
\usepackage[scr=rsfs,cal=euler]{mathalfa}
\usepackage{xfrac}
\usepackage[unicode,hidelinks,bookmarks=false]{hyperref}
\newcommand{\ie}{i.e.\ }
\newcommand{\cf}{cf.\ }
\newcommand{\resp}{respectively\ }
\newcommand{\Subsection}{\subsection}
\providecommand{\nobreakdash}{\nobreak\hbox{-}}
\setlength{\emergencystretch}{2em}
\multlinegap0pt
\usepackage{amssymb}
\usepackage{enumerate}
\newtheorem{thm}{Theorem}
 \def\BB{{\mathbb B}}
 \def\Jb{{\mathbf J}}
 \def\UU{{\mathbb U}}
 \def\Wb{{\mathbf W}}
 \def\d{\mathrm{d}}
\def\f{\frac}
 \def\fD{{\mathfrak D}}
 \def\fR{{\mathfrak R}}
 \def\g{{\gamma}}
 \def\l{{\lambda}}
 \def\la{{\langle}}
 \def\ra{{\rangle}}
\title{Bernstein Inequality on Parabolic Domains}
\author{Yuan Xu}
\date{Source: arXiv:2604.04268v1 (CC BY 4.0)}
\begin{document}
\maketitle

\noindent\emph{Source and licence.} Bernstein Inequality on Parabolic Domains. Version arXiv:2604.04268v1 (CC BY 4.0), distributed by arXiv under the Creative Commons Attribution 4.0 International licence, http://creativecommons.org/licenses/by/4.0/, as recorded in the arXiv licence metadata for this exact version and shown on the abstract page https://arxiv.org/abs/2604.04268v1. Full licence notice: \texttt{licenses/arxiv-2604.04268-CC-BY-4.0.txt}.\par\smallskip

\noindent\textit{Editorial scope.} Counted unit: the article's thm eq:fD\_UU2 (source0.tex lines 557--579). The article's complete original proof of it is delivered with it (source0.tex lines 581--633, one \texttt{proof} environment of 53 lines). No mathematics is written, repaired or re-derived: the statement and the proof are the article's own lines, in the article's own notation, and no retained line is substituted or re-flowed.  Retained uncounted as the source-local context the proof needs: lines 349--353; lines 464--467; lines 511--516.  Nothing was omitted from any retained span, so the package carries no omission record.  No retained line contains a citation, so the wrapper carries no bibliography and no bibliography entry.  No retained line contains a figure environment, an image-inclusion command or drawing code, so no image is delivered and the package declares no asset dependency.  Editorial formatting only, in addition: an AMS \texttt{amsart} preamble that restates the article's own declarations and macros used by the retained lines, namely the environments \texttt{thm} on separate counters, together with the standard packages those lines need.  The retained environments print in this document as Theorem 1 in order of appearance, whereas the article prints the same items with the numbering of its own full document.  The complete native source of the article as distributed in the arXiv e-print for this exact version is bundled unchanged as \texttt{sources/197873/source0.tex}.
\begin{equation} \label{eq:fD_Bd}
  \fD_\BB^\mu:= \Delta - \la x, \nabla \ra^2 - (2\mu + d-1) \la x, \nabla \ra 
  %= \sum_{i=1}^d (1-x_i^2) \partial_i^2 - 2 \sum_{1 \le i < j \le d} x_i x_j \partial_i \partial_j - (d+2\mu+1)
  %  \sum_{i=1}^d x_i \partial_i 
\end{equation}

\begin{equation}\label{intUU=}
  \int_{\UU^{d+1} } f(x,t) \d x \d t  = \int_0^b \int_{\|x\|^2 \le t} f(x,t) \d x \d t =
    \int_0^b t^{\frac{d}{2}} \int_{\BB^d} f\left(\sqrt{t} y, t\right) \d y \d t. 
\end{equation}

\begin{align}\label{eq:diff-eqnUU}
 & \fD_{\Jb}^{\g,\mu} u : = \left[ t (1-t) \partial_{tt} + (1-t) \la x,\nabla_x\ra \partial_t + \frac{1}{4} (1-t)\Delta_x  \right. \\
      & \qquad   \qquad  \quad
      +\left.   \left(\mu+\tfrac{d+1}{2}\right)(1-t)\partial_t   - \frac{\g+1}{2} (2t \partial_t +  \la x,\nabla_x\ra) \right]  u 
       = - \l_{m,n} u \notag
\end{align}

\begin{thm} 
Let $\g > -1$ and $\mu > -\f12$. The operator $ \fD_{\Jb}^{\g,\mu}$ can be rewritten as
\begin{align} \label{eq:fD_UU2}
 \fD_{\Jb}^{\g,\mu} =  \fR_\Jb^{\g,\mu}  + \frac{1-t}{4t} \fD_\BB^{\mu, (y)}
\end{align}
where $\fD_\BB^{\mu, (y)}$ denotes the spectral operator $\fD_\BB^\mu$, defined in \eqref{eq:fD_Bd}, acting on 
the variable $y = \frac{x}{\sqrt{t}} \in \BB^d$, and $\fR_\Jb^{\g,\mu} $ is defined by 
\begin{align*}
 \fR_\Jb^{\g,\mu} =  \frac{1}{t^{\f d 2} \Wb_\UU^{\g,\mu}(x,t)} \left( \partial_t + \frac{1}{2t} \la x ,\nabla_x\ra \right) 
  \left[t^{\f{d}{2} +1} \Wb_\UU^{\g+1,\mu}(x,t) \left( \partial_t + \frac{1}{2t} \la x ,\nabla_x\ra \right) \right]. \notag
 \end{align*}
In particular, $\fD_\Jb^{\g,\mu}$ is self-adjoint on $L^2(\UU^{d+1}, \Wb_{\UU}^{\g,\mu})$. Furthermore, 
\begin{align}\label{eq:int_fDU1}
  &  - \int_{\UU^{d+1}} \fR_{\Jb}^{\g,\mu} f(x,t)\cdot g(x,t) \Wb_{\UU}^{\g,\mu}(x,t) \d x \d t \\
 =  \int_{\UU^{d+1}} & t(1-t)  \left( \partial_t + \frac{1}{2t} \la x ,\nabla_x\ra \right) f(x,t)  \left( \partial_t + \frac{1}{2t} \la x ,\nabla_x\ra \right) g(x,t) \Wb_{\UU}^{\g,\mu}(x,t) \d x \d t \notag
\end{align}
and
\begin{align}\label{eq:int_fDU2}
 - & \int_{\UU^{d+1}}  \frac{1-t}{4t} \fD_{\BB}^{\mu, (y)} f(x,t)\cdot g(x,t) \Wb_{\UU}^{\g,\mu}(x,t) \d x \d t \\
  & =  \int_0^1 \frac{1-t}{4t} \left[ \int_{\BB^d} \fD_\BB^{\mu, (y)} f\left(\sqrt{t} y, t\right ) \cdot g\left(\sqrt{t} y, t\right)
   \Wb_\BB^\mu (y) \d y\right]  t^{\f{d-1}{2} + \mu}(1-t)^\g \d t. \notag
\end{align}
\end{thm}

\begin{proof}
The main hurdle lies in recognizing the correct form. The verification comes down to heavy computation of
derivatives, tedious but not difficult, and our proof will be succinct. We start from an observation 
$$
   \left(\partial_t + \frac{1}{2t} \la x,\nabla_x\ra \right) (t-\|x\|^2)^{\mu - \frac12}
    = \frac{\mu - \f12}{t} (t-\|x\|^2)^{\mu-\f12},
$$
which follows from a quick computation. This identity is then used to take the derivatives in $\fR_\Jb^{\g,\mu}$ to
deduce, after simplification, 
\begin{align*}
 \fR_\Jb^{\g,\mu} 
  = \,& \left[ \mu+ \frac {d+1} 2 - \left(\mu+\g+ \frac {d+3} 2 \right)t \right]  \left( \partial_t + \frac{1}{2t} \la x ,\nabla_x\ra \right) \\
    & \,\, + t(1-t)  \left( \partial_t + \frac{1}{2t} \la x ,\nabla_x\ra \right)^2, 
\end{align*}
Furthermore, the second term on the right-hand side satisfies 
\begin{equation*}
   t \left( \partial_t + \frac{1}{2t} \la x ,\nabla_x\ra \right)^2 = t \partial_{tt} + \la x, \nabla_x \ra \partial_t + 
      \frac{1}{2t} \left( \frac 12 \la x, \nabla_x\ra^2 - \la x, \nabla_x\ra \right),
\end{equation*}
so that we can deduce, after rearranging terms and using the definition of $\fD_\UU^{\g,\mu}$ in \eqref{eq:diff-eqnUU}, that 
\begin{align*}
 \fR_\Jb^{\g,\mu}& =  \fD_\UU^{\g,\mu} -\frac{1-t}4 \Delta_x + \left (\mu+\frac{d+1}{2}\right )\frac{1-t}{2t} \la x, \nabla_x \ra
  +   \frac{1-t}{2t} \left( \frac 12 \la x, \nabla_x\ra^2 - \la x, \nabla_x\ra \right)\\
  & =  \fD_\UU^{\g,\mu} - \frac{1-t}{4t} \left(t \Delta_x - \la x, \nabla_x \ra^2 - (2\mu+d-1) \la x, \nabla_x \ra \right).
\end{align*}  
Now, setting $y = \frac{x}{\sqrt{n}}$, it follows that $\frac{\partial}{\partial y_i} = \sqrt{t}  \frac{\partial}{\partial x_i}$
and $\la x, \partial_x\ra = \la y, \partial_y\ra$, so that 
$$
t \Delta_x - \la x, \nabla_x \ra^2 - (2\mu+d-1) \la x, \nabla \ra = \Delta_y - \la y, \nabla_y \ra^2 - (2\mu+d-1) \la y, \nabla_y \ra
 = \fD_\BB^{\mu, (y)}.
$$
Together, these identities prove \eqref{eq:fD_UU2}. 

Let $I_\fR$ be the integral on the left-hand side of \eqref{eq:int_fDU1}. Parameterizing the integral with $x = \sqrt{t} y$
as in \eqref{intUU=}, and using the identity
\begin{equation}\label{eq:dt_UU}
   \frac{\d}{\d t} f\left(\sqrt{t} y, t\right ) = \left(\partial_t + \frac{1}{2t} \la x, \nabla_x \ra \right) f\left(\sqrt{t} y, t\right )
\end{equation}
and $\Wb_\UU^{\g,\mu}(y, 1) = 0$, we deduce via integration by parts, 
\begin{align*}
  I_{\fR}\,& = \int_{\BB^d} \int_0^1 \frac{\d} {\d t} \left[ t^{\frac{d}{2}+1} \Wb_\UU^{\g+1,\mu}\left(\sqrt{t} y, t\right) 
     \frac{\d} {\d t}f\left(\sqrt{t} y, t\right ) \right] g \left(\sqrt{t} y, t\right ) \d t \d y \\
     & = - \int_{\BB^d} \int_0^1\left[ t^{\frac{d}{2}+1} \Wb_\UU^{\g+1,\mu}\left(\sqrt{t} y, t\right) 
   \frac{\d} {\d t}  f\left(\sqrt{t} y, t\right ) \right]   \frac{\d} {\d t} g \left(\sqrt{t} y, t\right ) \d t \d y,
% & = - \int_{\BB^d} \int_0^1\left[ t^{\frac{d}{2}+1} \Wb_{\g+1,\mu}\left(\sqrt{t} y, t\right) 
%   \frac{\d} {\d t}  f\left(\sqrt{t} y, t\right ) \right]   \frac{\d} {\d t} g \left(\sqrt{t} y, t\right ) \d t \d y   
\end{align*} 
which is equal to the right-hand side of \eqref{eq:int_fDU1} upon using \eqref{eq:dt_UU}, and \eqref{intUU=}. 

Finally, the identity \eqref{eq:int_fDU2} is an immediate consequence of rewriting the integral via $x = \sqrt{t} y$
and using $\Wb_\UU^{\g,\mu} (\sqrt{t} y, t) = (1-t)^\g t^{\mu-\f12} W_\BB^\mu (y)$. Since $\fD_\BB^\mu$ is self-adjoint,
the self-adjointness of $\fD_\Jb^{\g,\mu}$ follows from the two identites \eqref{eq:int_fDU1} and \eqref{eq:int_fDU2}. 
\end{proof}

\end{document}
